Let $n\ge 2$ and let $\pi(n)$ be the number of primes $\le n$. Fix an integer $k$ with $\pi(n)<k\le n$.
For a subset $A\subseteq\{1,\dots,n\}$ of size $|A|=k$, consider the following property for a given integer $r\ge 1$:
there exist distinct primes $p_1,\dots,p_r$ such that at least $r$ elements $a\in A$ have all their prime divisors contained in $\{p_1,\dots,p_r\}$.
Define $f(k,n)$ to be the smallest integer $r$ such that this property holds for every $A\subseteq\{1,\dots,n\}$ with $|A|=k$.

Ambiguity note (literal vs corrected):
- Literally, the condition ``$a$ is only divisible by primes from $\{p_1,\dots,p_r\}$'' is vacuously true for $a=1$ (since $1$ has no prime divisors). This makes very small-$n$ cases somewhat degenerate.
- A minimal corrected variant sometimes used in number theory is to require the counted elements $a\in A$ to satisfy $a>1$.
In the computations below, I report both conventions for small $n$.

The question asks whether
\[
2\pi(\sqrt n)-f(\pi(n)+1,n)\to\infty\quad\text{as }n\to\infty,
\]
and more generally to estimate $f(k,n)$ when $\pi(n)+1<k=o(n)$.
